%% kapa: exam=2 round=22 year=2011 period=2 subject=appraisal_theory subject_ko=감정평가이론 seq=3 number=3 points=20
%% source: Q-Net 감정평가사 기출문제 (cst00309) artlSeq=5040389 file=extracted/1008/5040389/2011년 제22회 감정평가사 2차시험 기출문제/2교시.hwp
%% conversion: kordoc 4.24.51 -> tools/gen2.py (auto); review=proofread-sonnet; flags=none
\begin{longquestion}{【문제 3】}
정비사업은 도시환경을 개선하고 주거생활의 질을 높이는 것이 목적인데 그 중 주택재개발사업은 정비기반시설이 열악하고 노후ㆍ불량건축물이 밀집한 지역의 주거환경을 개선하기 위한 사업이다. 이에 관한 감정평가사의 역할이 중요한 바, 다음의 물음에 답하시오. (20점) \par
\medskip\noindent (1) 주택재개발사업의 추진 단계별 목적에 따른 감정평가업무를 분류하고 설명하시오. (10점) \par
\medskip\noindent (2) 종전자산(종전의 토지 또는 건축물)과 종후자산(분양예정인 대지 또는 건축물의 추산액)과의 관계를 설명하시오. (10점) \par
\end{longquestion}
